% F01：本书原创矢量示意；依据所在节的定义、证明或明确模型。
\begin{figure}[htbp]
\centering
\begin{tikzpicture}[x=1cm,y=1cm,>=stealth,every node/.style={font=\small},line width=.65pt]

\fill[AnalysisBlue!12] (0,0) rectangle (2.3,2.3);
\draw[thick,AnalysisBlue] (0,0) rectangle (2.3,2.3);
\fill[AnalysisBlue!12] (7,0)--(9.3,0)--(10.5,2.3)--(8.2,2.3)--cycle;
\draw[thick,AnalysisBlue] (7,0)--(9.3,0)--(10.5,2.3)--(8.2,2.3)--cycle;
\draw[orange!80!black,very thick] (0,1)--(2.3,1);
\draw[orange!80!black,very thick] (7.522,1)--(9.822,1);
\draw[->] (3.2,1.2)--(5.9,1.2) node[midway,above] {\(S_t(x,y)=(x+ty,y)\)};
\node at (1.15,-.45) {原矩形};\node at (8.8,-.45) {剪切后的像};
\node[orange!80!black] at (5.3,2.6) {同一高度的截面等长};

\end{tikzpicture}
\caption[剪切为什么保持体积]{剪切为什么保持体积。剪切保留水平截面的长度；由非负截面积分得到面积不变。图示为正剪切参数的一个情形。}
\label{b1:fig:visual-f01}
\end{figure}
